\section{剩余的比较为什么困难}\label{sec:difficulty}
下界证明将每类方向的截面下界积分到互不相交的径向区域，由同一个实际面积控制总和。构造之所以成功，是因为它们安排不同记录反复覆盖同一份材料。要弥合剩余间隙，这两种描述必须彼此约束：高效重叠所依赖的相位对齐，必须提供能用于每个输入的下界信息。

第一个障碍是对完整记录的有损压缩。纵向分拆同时固定两个端点相位，有符号高度决定这些端点如何与实际射线相遇，而底边的每个内部点都在中间半径处参与其中。只保留端点分布、各半径上的占据长度或孤立圆周采样，会丢掉实际角度之间的对应关系。证明这类简化模型的最优值，并不自动解决原始完整单位线段问题。

完整精确的径向函数并没有这种信息损失：它确定整个紧原点星体 $K$。例如，给定一条候选单位段 $[a,b]$，置 $x_t=(1-t)a+tb$，就有
\[
 [a,b]\subset K\quad\Longleftrightarrow\quad
 \forall t\in[0,1],\quad x_t\ne0\ \Longrightarrow
 \|x_t\|\le\rho_K\!\left(\frac{x_t}{\|x_t\|}\right).
\]
只知道每个半径处的占据总弧长，就不能作这种逐点检验。若只保留径向函数的几乎处处等价类，还须恢复全方向的实际代表元；第~\ref{sec:strong-metric} 节已经完成这一步，并证明相应完备化的面积下确界仍为 $A_*$。困难在于建立能控制全部这些完整段的锐面积不等式，而不是精确径向表示丢掉了集合。

第二个障碍是净面积比较。连接相邻方向块、旋转或置换方向块，以及压缩高度，都可以保持单位线段可行性。但桥接或替换仍可能在保留背景之外新增材料。在面积符号尚未解决的路线中，缺失的步骤是用实际撤去的材料控制新增并集，而不是再次验证单位线段公式。式~\eqref{eq:physical-change} 指出了缺失的不等式。双剪切例子更进一步：对它的固定规则，即使任意接近下确界，所希望的符号也不成立。

第三个障碍涉及量词。一个固定输入可以有严格的局部改进，但改进量沿复杂度不断增加的序列趋于零。对固定单元数有效的归约，在单元数变化时可能失去误差控制。在一个周期族内最优的正规形，可能无法表示任意有限扫掠。给受限族的未知下确界换一个名字，解决不了这些问题。一个保值归约必须构造合法的完整单位针，核算撤去的材料，并对它所应用的近极小化序列保持一致的面积误差控制。

连续性带来另一种约束。短轴障碍证明，几乎保留一个固定体的材料，可能使任何连续单位针截面都无法存在，即使不限制半径。它不排除代价更低的自由重构。因此，$A_{\mathrm{coh}}=A_*$ 仍是一个有意义的问题，但不能靠具有连续闭单位针截面的体的强稠密性来作一般证明。论证中的零代价径向状态也解释了逐点配置惩罚为什么不够：证明还需要方向窗口上的集体面积下界来排除它们。

最后，内禀最长弦的界避开了来源标签的选择，却引入了自身的损失。联合壳层估计只看到一条长完整弦的短端部，再借助极大函数晕区恢复其来源方向。这正确支付了所有尺度，却未能足够精细地区分有竞争力构造的重叠模式。丢掉这些区别后再改进一个标量，不能代替找回缺失的几何信息。

因此，无须再引入一个辅助阈值，就能陈述剩余问题。我们需要一个对所有类别成立的面积不等式，其值能由合法有限族逼近；或者需要一个真正控制面积的归约，把近极小者送入一个值可以确定的类别。$1/10$ 定理提供一致的基线。严格的上界改进说明，已知构造中仍有可利用的材料。反例指出了哪些归约不能把这些构造转移到任意输入。这些结果共同指出，从这些截面估计走到真正的锐值定理，还须完成哪些工作。
